\documentclass[11pt,a4paper]{article}
\usepackage[T1]{fontenc}
\usepackage[utf8]{inputenc}
\usepackage{lmodern,microtype}
\usepackage[a4paper,top=22mm,bottom=22mm,left=25mm,right=25mm,headheight=14pt]{geometry}
\usepackage{amsmath,amssymb,amsthm,mathtools,booktabs,array,tabularx,xcolor,fancyhdr,enumitem}
\usepackage[hyphens]{url}
\usepackage[colorlinks=true,linkcolor=black!78!blue,urlcolor=blue!45!black,citecolor=blue!45!black]{hyperref}
\definecolor{ink}{HTML}{213849}
\definecolor{accent}{HTML}{16645E}
\pagestyle{fancy}\fancyhf{}\fancyhead[L]{\small\sffamily\color{ink} Catalan determinants and binomial convolutions}
\fancyhead[R]{\small\sffamily\color{ink} Keio research manuscript}\fancyfoot[C]{\small\thepage}
\renewcommand{\headrulewidth}{0.3pt}
\newtheorem{theorem}{Theorem}[section]\newtheorem{lemma}[theorem]{Lemma}
\newtheorem{proposition}[theorem]{Proposition}\newtheorem{corollary}[theorem]{Corollary}
\theoremstyle{definition}\newtheorem{problem}[theorem]{Auxiliary problem}
\theoremstyle{remark}\newtheorem*{remark}{Remark}
\newcommand{\Z}{\mathbb Z}
\newcommand{\Q}{\mathbb Q}
\newcommand{\qbin}[2]{\genfrac{[}{]}{0pt}{}{#1}{#2}_{\!q}}
\setlist{topsep=4pt,itemsep=2pt}
\allowdisplaybreaks[1]\setlength{\emergencystretch}{2em}
\hypersetup{pdftitle={Exact degrees of Catalan Hankel determinants and congruences for binomial convolutions},pdfauthor={Keio and Shaitetsu}}
\title{\color{ink}\bfseries Exact degrees of Catalan Hankel determinants\\[3pt] and congruences for binomial convolutions}
\author{Keio and Shaitetsu}
\date{Research manuscript \enspace$\cdot$\enspace 20 November 2026}
\begin{document}
\maketitle
\begin{abstract}
We prove the odd-power degree conjecture of Wang and Xin for all Catalan convolution powers in its stated range. For arbitrary integral shifts, the signed determinant on each prescribed residue class is a polynomial, as proposed by Cigler. We also determine the even unshifted degrees and the degree and reciprocal sign of the even generating-function numerator. Equality of degrees follows from nonvanishing, since the leading minor in a confluent Chebyshev determinant is a Gram determinant of an explicitly positive discrete measure. In a second application of finite reflection, we prove Zhi-Wei Sun's Conjecture 35(i) for every modulus, both signs and arbitrary powers. A two-index reflection on the nonzero digits of a Gaussian-binomial convolution gives a polynomial divisibility theorem with independently variable row weights. We identify the sole exception in Cigler's odd numerator degree formula. For Sun's related cubic-weight problems we prove the specified prime-square congruences, including the exceptional residue evaluations.
\end{abstract}

\section{Statements and provenance}
Let $C(z)=1+zC(z)^2$ with $C(0)=1$. Extend the coefficients of $C(z)^r$ by zero in negative degrees. We study their contiguous, shifted Hankel determinants and an integer-valued polynomial built from two complementary binomial rows. These are two independent applications of reflection. In the determinant, folding frequencies isolates a minor whose leading coefficient must survive. In the congruence, reflection pairs terms of opposite weight at each relevant root of unity.

\begin{center}
\small\renewcommand{\arraystretch}{1.15}
\begin{tabularx}{\linewidth}{@{}p{0.34\linewidth}X@{}}
\toprule
Original target & What is proved here\\\midrule
Wang--Xin 9 \cite{WX} & Exact degrees for both odd-power residue classes over the entire range $1\le j\le t$.\\
Cigler 4, 6, 12 \cite{Cig} & Polynomiality at every integral shift, and both published unshifted degree formulas. Cigler 12 is included in Wang--Xin 9.\\
Cigler 15 \cite{Cig} & The specified even denominator clears the series. Its numerator has the exact predicted degree and reciprocal sign.\\
Cigler 16 \cite{Cig} & The previously known exception $(k,m)=(1,0)$ is the sole degree exception. All proposed reciprocal signs hold.\\
Cigler 18 \cite{Cig} & A four-value exact certificate refutes the printed odd leading coefficient.\\
Sun 35(i) \cite{Sun} & Integer-valuedness for every positive $\ell,m,n$, every integer argument and each sign.\\
\bottomrule
\end{tabularx}
\end{center}

Wang and Xin proved Theorem~11 for $r\le31$ and reported further computations for Conjecture~9 through $r=41$. Their printed Conjecture~10 quantifies the asymmetric expression $(2j-1)(t-j)$ on both residues for every $1\le j\le t$. It gives incompatible degrees at reflected indices. For example, $t=2$ gives one degree at $j=1$ and degree zero at $j=2$ for the same pair. The even result proved here is Cigler's separately published, symmetric Conjecture~6, equation~(19). The odd result uses Wang and Xin's original quantifiers. Cigler's statements are taken from version~2 of \cite{Cig}. His equations (22)--(27) already connect cotangent and tangent coefficients with particular residue classes. We identify the determinant of such coefficients in the general degree problem and prove its strict nonvanishing.

Krattenthaler \cite{Krat} gives the general finite determinant reduction. Fulmek and Cigler \cite{Back} give the backward-shift identity used below. Chern--Shi and Han \cite{Other} evaluate related families. The public proof of Cigler 11 \cite{Omega} already combines Catalan moments, a polynomial basis, and finite-dimensional reduction in a particular residue class. The remaining difficulty for our degree formulas is the nonvanishing of the leading coefficient simultaneously in every prescribed residue.

The motivating corpus snapshot \cite{Corpus} has 217 numbered TeX manuscripts. Two particularly relevant constructions are the finite Christoffel reduction of Hankel determinants in manuscript 079 and the positive moment expansion of a leading determinant term in manuscript 026. The two are combined here at the multiple zero of a moment multiplier. Manuscript 208 studies cyclotomic vanishing of a different double sum. Our arithmetic argument constructs complementary digit rectangles at every divisor of the modulus. None of the proofs below requires accepting a new theorem from the corpus.

\paragraph{The auxiliary determinant problem.}
Suppose column operations on a polynomial evaluation determinant leave a block of low Taylor coefficients at a multiple root. An entrywise degree bound gives only an upper bound on the block determinant. When does its largest possible degree actually occur? The following criterion gives the noncancellation step used in both parities.

\begin{lemma}\label{lem:leading}
Fix $U\ge V\ge1$. Suppose that for the integers $h$ between $U-V+1$ and $U+V-1$,
\[
 a_h(H)=H^{2h-1}\{b_h+O(H^{-1})\},\qquad
 b_h=\lambda\int_{(0,\infty)}s^{h}\,d\mu(s)
\]
as $H\to\infty$, where $\lambda\ne0$, $\mu$ is positive with infinite support, and all indicated moments are finite. Then $\det(a_{U+i-j}(H))_{0\le i,j<V}$ has a nonzero leading term of order $H^{2UV-V}$.
\end{lemma}
\begin{proof}
The sum of the indices $U+i-\sigma(i)$ is $UV$ in every permutation. The leading coefficient is therefore the determinant of the $b$'s. Reversing its columns gives, apart from $\lambda^V$ and the reversal sign, the Gram matrix of $1,s,\dots,s^{V-1}$ for the measure $s^{U-V+1}d\mu(s)$. A nonzero polynomial has finitely many roots, so its squared norm is strictly positive.
\end{proof}


\section{Polynomiality for every integral shift}
Write $C(z)=1+zC(z)^2$, $C(0)=1$, and set
\[
c_{r,u}=[z^u]C(z)^r,\quad
D_{r,m}(N)=\det(c_{r,i+j+m})_{0\le i,j<N}.
\]
Here $c_{r,u}=0$ for $u<0$ and $D_{r,m}(0)=1$. Indices of determinants are nonnegative unless explicitly extended as functions. Cigler \cite[Conj.~4, eqs.~(16)--(17)]{Cig} asks for polynomiality on the following residue classes. \begin{theorem}[Cigler 4]\label{result:allshift} For $t\ge1$ and every integer $m$,
\[
(-1)^{n\binom t2}D_{2t,m}(tn+j)\quad(0\le j<t),\qquad
(-1)^{tn}D_{2t+1,m}((2t+1)n+j)\quad(0\le j\le2t)
\]
agree with rational polynomials in $n$ for every integer $n\ge0$.
\end{theorem}
We prove both assertions, including the initial zero entries of arbitrarily deep backward shifts.

Let $S_{-1}(y)=0$, $S_0(y)=1$, $S_{u+1}(y)=yS_u(y)-S_{u-1}(y)$ and extend the recurrence in both directions.
We use two monic polynomial bases
\[
E_u(x)=S_u(x-2),\qquad O_u(x)=S_u(x-2)+S_{u-1}(x-2)\qquad (u\ge0).
\]
In each case construct a tridiagonal matrix $J$ on indices $0,1,\ldots$ with both off-diagonals equal to $1$ and diagonal $2$, except that $J_{00}=1$ for $O$. Put $L(f)=(f(J))_{00}$. The chosen basis $B=E$ or $O$ satisfies $B_i(J)e_0=e_i$ by the recurrence, hence $L(B_i B_j)=\delta_{ij}$. The series $G_b(z)=\sum_{u\ge0}L(x^u B_b)z^u$ is
\[
 G_b(z)=z^b C(z)^{2b+2}\quad(B=E),\qquad
 G_b(z)=z^b C(z)^{2b+1}\quad(B=O).
\]
Indeed the generating vector satisfies $G_0=1+z(J_{00}G_0+G_1)$ and $G_b=z(G_{b-1}+2G_b+G_{b+1})$ for $b>0$; substitution using $C=1+zC^2$ verifies these equations and the coefficients are unique recursively. Consequently the shifted sequence is exactly $L(x^u q)$ for $u\ge0$ in the following cases:
\[
\begin{array}{c|cccc}
r&B&b&a&q(x)\\\hline
2t&E&t-1&m+t-1&x^a E_b(x)\\
2t+1&O&t&m+t&x^a O_b(x).
\end{array}
\]
Here we first assume $a\ge0$; this dictionary also includes negative $m$ if the displayed exponent permits it.

Put $d=a+b$. For any $f$, let $V(f)$ be the vector of its Taylor coefficients $[x^i]f$ for $0\le i<a$, followed by its values at the $b$ distinct roots of $B_b$. These roots are nonzero, so the mapping on polynomials of degree less than $d$ is invertible by elementary polynomial division. Define
\[
F(N)=(-1)^{Nd}
\frac{\det[\,V(B_{N+s})\,]_{s=0}^{d-1}}{\det[\,V(B_s)\,]_{s=0}^{d-1}}
\qquad (N\in\mathbb Z).
\]
The displayed vectors are columns; empty determinants are $1$. For $N\ge0$, $F(N)=D_{r,m}(N)$. To see this, on the polynomial space of degree $<N+d$ take the coordinate map comprising the $N$ functionals $L(B_i f)$, $0\le i<N$, followed by $V(f)$. In the ordered monic basis $B_0,\ldots,B_{N+d-1}$ its determinant is the numerator determinant without sign. In the ordered basis $1,x,\ldots,x^{d-1},q,xq,\ldots,x^{N-1}q$, which is also monic by degree, its determinant is $(-1)^{Nd}\det(V(B_s))_{s<d}\,D_{r,m}(N)$. This follows by expanding on the last $d$ rows and replacing each $B_i$ in the top block by $x^i$ by unitriangular row operations.

For even powers the roots are $2+2\cos(\ell\pi/t)$, $1\le\ell<t$, and $B_u$ there equals $\sin((u+1)\ell\pi/t)/\sin(\ell\pi/t)$. For odd powers the roots are $2+2\cos(2\ell\pi/(2t+1))$, $1\le\ell\le t$. Indeed $O_u(2+2\cos\theta)=\sin((u+1/2)\theta)/\sin(\theta/2)$. These formulas hold for all integer $u$. For $i\ge0$,
\[
 [x^i]E_u(x)=(-1)^{u-i}\binom{u+i+1}{2i+1},
 \qquad [x^i]O_u(x)=(-1)^{u-i}\binom{u+i}{2i}.
\]
The first formula follows by expanding $(1-(x-2)z+z^2)^{-1}$, the second by Pascal's identity. The polynomial-binomial interpretation extends them to negative $u$ by the backwards recurrence (or $E_{-u-2}=-E_u$, $O_{-u-1}=-O_u$).

Take $N=pn+j$ with $p=t$ or $2t+1$ respectively and any fixed integer $j$. Increasing $n$ by $1$ multiplies each simple-root row by $(-1)^\ell$ in the even case and leaves it unchanged in the odd case. Every jet row has a common factor $(-1)^{pn}$ times a row of polynomials in $n$. With the exterior determinant sign included, the net sign per step is $(-1)^{t(d+a)+\binom t2}=(-1)^{\binom t2}$ in the even case and $(-1)^{(2t+1)(d+a)}=(-1)^t$ in the odd case. Thus $F(pn+j)$ has the desired form on all integer $n$. Its polynomial has rational coefficients, since at all sufficiently large $n$ it takes integer values after extracting a sign, and interpolation at integer arguments is over $\mathbb Q$.

It remains to extend this argument to $a<0$. Write $m=-h$. The backward-shift identities of Fulmek and Cigler (Cigler, \cite{Back}, Theorems~1--2) give
\[
D_{r,-h}(N)=(-1)^{\binom{h+1}2}D_{r,h-r+2}(N-h-1)\quad(N\ge h+1).
\]
For $1\le N\le h$ the first matrix row vanishes. In the reflected determinant the multiplier degree is $d=h$ for even powers and $d=h+1$ for odd powers, with nonnegative zero multiplicity in both cases. Use its extended $F$ just defined. For $1\le u\le h$ the columns of $F(-u)$ contain $E_{-1}=0$ in the even case and both $O_{-1}=-O_0$ and $O_0$ in the odd case, so $F(-u)=0$. At $N'=-h-1$, the even columns are $-E_{h-1},\ldots,-E_0$, and the odd columns are $-O_h,\ldots,-O_0$. In the even case the extension sign $(-1)^{N'd}$ equals $1$, while reversal and column signs contribute $(-1)^{h+\binom h2}$. In the odd case $d=h+1$ and the total sign is $(-1)^{d^2+d+\binom d2}$. Both give $F(-h-1)=(-1)^{\binom{h+1}2}$. Hence the extended reflected formula also gives the actual value $1$ at $N=0$ and every initial zero. This proves Conjecture~4.

\section{A noncancellation theorem for the exact degrees}
The numerical degrees in the even, unshifted case were investigated by Wang and Xin (Discrete Math.\ 342 (2019), 2694--2716). We prove the symmetric half-range statement in Cigler, arXiv:2308.07642v2, Conjecture~6, equation (19). For $t\ge1$ express the residues by $1\le j\le t$, with $j=t$ denoting the translated residue zero. Write $A=t-j$ and $B=j-1$.\begin{theorem}[Cigler 6] The degree of the normalized polynomial associated to $D_{2t,0}(tn+j)$ is exactly
\[
\boxed{\ (2\min(A,B)+1)\max(A,B).\ }
\]
\end{theorem} In particular the degree for $1\le j\le(t+1)/2$ is $(2j-1)(t-j)$ and the residue $t+1-j$ has the same degree. Translating the variable by $1$ at residue zero preserves degree. For $t=1$ all determinants equal one.

Here $a=b=t-1$, $d=2t-2$, and $N=tn+j$. Write $P_u=E_{u-1}$, $M=t(n+1)$ and $K=t(n+2)$. In the evaluation determinant the frequency indices $u$ run through consecutive integers
\[
M-A+1,\ldots,K+B-1.
\]
There are $A$ kernel columns centered at $M$: when $A>0$, take $P_M$ and the sums $P_{M+i}+P_{M-i}$ for $1\le i<A$. Similarly there are $B$ kernel columns at $K$. The identity
\[
P_{H+i}+P_{H-i}= P_H\,(S_i(x-2)-S_{i-2}(x-2))\quad(i\ge1)
\]
follows from sine addition; the parenthesized factors are monic of degree $i$. The two groups are disjoint since $A+B=t-1$. All these columns vanish in the simple-root rows. In each pair keep one original column and replace the other by the sum. Keep the unpaired columns as well. These operations and fixed permutations have constant nonzero determinant. The simple-root matrix has rank $t-1$: the original interval has length $2t-2$ and contains the entire block $M+1,\ldots,M+t-1$ (also when $A=0$ or $B=0$, with the corresponding endpoint included). Its values form a sine transform matrix, invertible by orthogonality, up to row signs. Thus the $t-1$ complementary columns give a nonsingular simple-root block, whose determinant has absolute value independent of $n$. Expanding on this block, and replacing the monic factors in each group by monomials, reduces the absolute value, up to a constant independent of $n$, to the determinant of the first $a=A+B$ Taylor coefficients of
\[
x^i P_M\quad(0\le i<A),\qquad x^i P_K\quad(0\le i<B).
\]

Write $U=\max(A,B)$ and $V=\min(A,B)$. Let $H$ be the frequency of the larger group and $H+\delta$ that of the smaller group, so $\delta=t$ or $-t$. Multiplication by $P_H$ in $\mathbb R[x]/(x^a)$ has determinant $P_H(0)^a$, with $|P_H(0)|=H$. Dividing all columns by $P_H$ and expanding along $1,x,\ldots,x^{U-1}$ leaves
\[
 H^a\left|\det\bigl([x^{U+i-j}]R_H(x)\bigr)_{0\le i,j<V}\right|,
 \qquad R_H(x)=P_{H+\delta}(x)/P_H(x).
\]
When $V=0$ the determinant is $1$. Otherwise every index $h=U+i-j$ in it is positive. For fixed $h\ge1$,
\[
 [x^h]\,R_H(x)/R_H(0)
 =\delta\,H^{2h-1}\left([z^h](\sqrt z\cot\sqrt z-1)+O(H^{-1})\right).
\]
For proof put $x=z/H^2$, $\phi=2\arcsin(\sqrt z/(2H))$. Then $P_u(z/H^2)=(-1)^{u-1}\sin(u\phi)/\sin\phi$. Dividing the quotient by its value at zero gives
\[
 \frac{H}{H+\delta}\left(\cos(\delta\phi)+\cot(H\phi)\sin(\delta\phi)\right)
 =1+\frac{\delta}{H}(\sqrt z\cot\sqrt z-1)+O(H^{-2}).
\]
The expansion is analytic and uniform for $|\sqrt z|$ on a fixed small circle, so Cauchy's formula gives the same estimate coefficientwise. Also $|R_H(0)|=(H+\delta)/H\to1$.

Reversing the columns of the leading coefficient matrix gives, up to common signs, the Hankel matrix $(c_{U-V+1+i+j})_{0\le i,j<V}$ with
\[
 c_h=2\pi^{-2h}\zeta(2h)\quad(h\ge1).
\]
This follows from the logarithmic derivative of Euler's product for $\sin y/y$. For $w=U-V+1\ge1$ its quadratic form at a nonzero polynomial $f$ of degree $<V$ equals
\[
 2\sum_{s\ge1}(\pi s)^{-2w} f((\pi s)^{-2})^2>0.
\]
The infinite sequence of distinct positive arguments prevents a nonzero polynomial from vanishing at all of them. Thus Lemma~\ref{lem:leading} applies and the leading determinant survives. Every permutation term has common scale $H^{2UV-V}$ because the sum of its coefficient indices is $UV$. Restoring $H^a$ gives the nonzero leading scale $H^{a+2UV-V}=H^{(2V+1)U}$. Polynomiality proved above turns this asymptotic statement into the asserted exact degree.

\section{Odd powers and the Wang--Xin odd degree conjecture}
Put $p=2t+1$. \begin{theorem}[Wang--Xin 9 and Cigler 12]\label{result:odd} For $1\le j\le t$ both normalized sequences for residues $j$ and $p+1-j$ are polynomials of degree
\[
 (j-1)(p-2j).
\]
\end{theorem} At $j=1$ the second residue is $p$, which translates residue zero. This establishes the whole statement of Wang and Xin, arXiv:1811.00248, Conjecture~9, and Cigler's Conjecture~12 for its indicated interior range.

Here $a=b=t$ and $d=2t=p-1$. Use the half-integral frequency $u=N+s+1/2$, so the column polynomial is $O_{u-1/2}$. When $x=2+2\cos\theta$ this column is $\sin(u\theta)/\sin(\theta/2)$. Set $A=t-j+1$ and $B=j-1$, so $A+B=t$. For $N=pn+j$, take the two centers
\[
 M=(n+1/2)p,\qquad K=(n+1)p .
\]
The frequency interval is $[M-(A-1),K+B-1/2]$ in steps of one. For $N=pn+p+1-j$ take instead $K=(n+1)p$ and $M=(n+3/2)p$; the interval in that case is $[K-B+1/2,M+(A-1)]$. Thus in either case take the singleton at $M$ and the $A-1$ pair sums at $M\pm i$ for $1\le i<A$, together with the $B$ pair sums at $K\pm(i+1/2)$ for $0\le i<B$. The groups are disjoint since $|M-K|=p/2$ and the two half-widths total less than $p/2$. Each of these $t$ columns vanishes at all the simple nodes $\theta=2\ell\pi/p$. The $p-1$ consecutive original frequencies occupy all but one of the $p$ residue classes of $u-1/2$. Their evaluations still span the $t$ simple rows: the classes with frequencies $u=1/2,3/2,\ldots,t-1/2$ and their negatives modulo $p$ come in pairs, each represented. Representatives of these pairs give an invertible evaluation matrix because $O_0,\ldots,O_{t-1}$ are monic of consecutive degrees. Keeping one original column per pair operation gives a complementary block. Its absolute determinant is independent of $n$, since shifting by $p$ fixes every simple-root evaluation.

For the zero jets put $x=2-2\cos\phi$. Since $u$ is half-integral,
\[
 O_{u-1/2}(x)=(-1)^{u-1/2}\frac{\cos(u\phi)}{\cos(\phi/2)}.
\]
The first group spans the columns $x^i U$, $0\le i<A$, and the second group the columns $x^i V$, $0\le i<B$, where
\[
 U=\frac{\cos(M\phi)}{\cos(\phi/2)},\qquad
 V=\sin(K\phi)\tan(\phi/2)=\frac{x}{2}S_{K-1}(2-x).
\]
Indeed, the pair sums at $M$ factor as $U$ times signed multiples of $\cos(i\phi)$, and those at $K$ factor as $V$ times signed multiples of
$\sin((i+1/2)\phi)/\sin(\phi/2)=1+2\sum_{v=1}^i\cos(v\phi)$. The displayed factors in each group are polynomials of successively increasing degrees with nonzero leading coefficients independent of $n$ in absolute value. The jet determinant hence changes only by a nonzero constant in absolute value. In particular $U(0)=1$ and $(V/x)(0)=K/2$.

If $A\ge B$, divide by $U$ in the truncated Taylor ring and expand on the monomials as before. The $B$-row determinant has entries $[x^{A+i-j'}](V/U)$, where $0\le i,j'<B$. For $x=z/M^2$ the trigonometric formulas, with $K/M=1+O(M^{-1})$, give
\[
 \frac{V}{U}(z/M^2)=\frac{1}{2M}\sqrt z\tan\sqrt z+O(M^{-2}).
\]
All estimates here and below hold uniformly on a sufficiently small fixed circle, so coefficient extraction is justified. The coefficients of $\sqrt z\tan\sqrt z$ at each $h\ge1$ are $2\sum_{s\ge1}(\pi(s-1/2))^{-2h}$, by the product for the cosine. Their shifted Hankel determinants of every order and shift at least one are strictly positive by the same quadratic-form argument as before. Thus the determinant has nonzero leading order $M^{2AB-B}$, including $B=0$ by the empty determinant convention.

If $B>A$, expand the original jet determinant along its constant row, where only the unshifted $U$ column survives. Every remaining column is divisible by $x$. Removing that factor turns the remaining jet matrix into one of size $t-1$, with $A-1$ columns generated by $U$ and $B$ columns generated by $V/x$. Divide by the latter generator, extracting the factor $(K/2)^{t-1}$. The remaining $(A-1)$-row minor now uses indices $h=B+i-j'\ge1$ of the ratio
\[
 \frac{U}{V/x}(z/M^2)
 =\frac{2}{M}\sqrt z\cot\sqrt z+O(M^{-2}).
\]
This follows directly from the displayed trigonometric formulas and is valid at $z=0$ by analytic continuation. The cotangent coefficient matrix has the strictly nonzero shifted Hankel determinant already proved. The resulting exponent is
\[
 (t-1)+2B(A-1)-(A-1)=B(2A-1).
\]
This includes $A=1$ with a zero-dimensional remaining determinant. Both ranges and both mirror residues therefore give degree $B(2A-1)$ as asserted.

\section{Even generating functions and their exact reciprocal signs}
Cigler's Conjecture~15, arXiv:2308.07642v2, equation (28), concerns every $t\ge1$ and $m\ge1-t$. In the even notation above put
\[
 a=m+t-1,\quad d=a+t-1,\quad L=1+\binom{a+1}{2},\quad
 \varepsilon=(-1)^{\binom t2},\quad
 G(z)=\sum_{N\ge0}D_{2t,m}(N)z^N.
\]
Then $P(z)=(1-\varepsilon z^t)^L G(z)$ is a polynomial of exact degree
\[
 T=t\binom{a+1}{2}-a=\frac{a(ta+t-2)}2,
\]
and obeys $z^T P(1/z)=\sigma P(z)$, where
\[
 \sigma=-(-\varepsilon)^L(-1)^{\binom{d+1}{2}} .
\]
The degree is Cigler's second $(t+2)$-gonal number at $a$. Explicitly the pairs $(t\bmod4,m\bmod4)$ giving $\sigma=+1$ are
\[
 (0,1),(0,3),\qquad(1,0),(1,1),(1,2),(1,3),\qquad
 (2,0),(2,3),\qquad(3,1),(3,2).
\]
All other pairs give $-1$. This is exactly the stated palindrome and skew-palindrome pattern.

For the proof, expand the determinant defining $F(N)$ on the simple-root rows. After signs have been extracted on a fixed residue class, each jet minor is a determinant $\det(f_i(N+s_j))_{0\le i,j<a}$, where $\deg f_i=2i+1$ and the $s_j$ are fixed. Taylor expansion $f_i(N+s_j)=\sum_{v\ge0}f_i^{(v)}(N)s_j^v/v!$ and Cauchy--Binet give a sum indexed by distinct nonnegative derivative orders. Their sum is at least $\binom a2$, so the minor degree is at most $a^2-\binom a2=\binom{a+1}2$. The empty minor satisfies the same bound. Hence each normalized residue polynomial has degree at most $L-1$, proving that $P$ is a polynomial by the finite-difference identity.

The bilateral determinant continuation has the symmetry
\[
 F(-d-1-N)=\eta F(N),\qquad \eta=(-1)^{\binom{d+1}{2}}.
\]
In fact $E_{-u-2}=-E_u$ reverses the order of all $d$ columns and changes their signs. The extra determinant prefactor contributes $(-1)^{(-d-1-2N)d}=1$. Also $F(-s)=0$ for $1\le s\le d$ since then a zero $E_{-1}$ column occurs. A bilateral polynomial times a periodic sign has the rational reciprocity identity
\[
 G(1/z)=-\sum_{s\ge1} F(-s)z^s
\]
as a rational function expanded at $z=0$. To verify this, decompose the sequence over $\mathbb C$ into a finite linear combination of $N^h\lambda^N$ by Fourier expansion of its period-$2t$ coefficients. The identity for $h=0$ is the geometric series; applying the Euler derivative $z\,d/dz$ proves every other case with the sign change at $1/z$ included. The vanished strip and the symmetry now give $G(1/z)=-\eta z^{d+1}G(z)$. The identity for $P$ follows with $T=tL-d-1$. Since $P(0)=1$, the leading coefficient at this stated degree equals $\sigma$, including the possibility $T=0$. Finally $\binom v2\bmod2$ depends only on $v\bmod4$; substitution in the displayed sign yields the table above.


\section{The complete degree and reciprocity classification for odd numerators}
Cigler's Conjecture~16 in arXiv:2308.07642v2, equation~(30), uses
$k\ge1$, $m\ge1-k$, and the multiplier
$(1+(-1)^k z^{2k-1})^{1+\binom{k+m}{2}}$.
The Omega Institute \cite{Omega} has given a public formal refutation of its exact
degree at $(k,m)=(1,0)$. Its dossier explicitly lists the corrected
degree and general reciprocity as open there. We determine the actual
degree and reciprocal sign for every allowed pair.

Put
\[
 p=2k-1,\qquad t=k-1,\qquad a=m+t,\qquad
 d=a+t,\qquad L=1+\binom{a+1}{2},\qquad \varepsilon=(-1)^t,
\]
and define
\[
 G(z)=\sum_{N\ge0}D_{p,m}(N)z^N,\qquad
 Q(z)=(1-\varepsilon z^p)^L G(z).
\]
Then $Q$ is always an integer polynomial with $Q(0)=1$.
For $(k,m)\ne(1,0)$ its exact degree and reciprocal law are
\[
 T=pL-d=\pi_{2k+1}(a)+k,\qquad
 z^TQ(1/z)=(-1)^{1+kL+\binom d2}Q(z),
 \qquad \pi_s(a)=\frac{a((s-2)a+s-4)}2.
\]
For $(k,m)=(1,0)$ one has $Q(z)=1$, of degree zero. Thus the known
exception is the only exception to the proposed degree. The sign claims
in the original hold throughout. More precisely, the reciprocal sign
is positive for $k\equiv1\pmod4$ with $m$ even, for
$k\equiv3\pmod4$ with $m$ odd, and for $k$ even with
$m\equiv0,1\pmod4$. It is negative in every other case.

\paragraph{Proof.}
The odd moment construction above applies also when $t=0$, by taking
$O_0=1$ and an empty simple-root block. Since $a\ge0$, the bilateral
determinant $F(N)$ is already defined and agrees with $D_{p,m}(N)$
for every $N\ge0$. Expand its determinant along the simple-root rows.
On each residue class, after the factor $\varepsilon^n$ is removed,
each Taylor-coefficient minor has the form
$\det(f_i(N+s_j))_{0\le i,j<a}$, up to constant signs, with
$\deg f_i=2i$ and fixed $s_j$.
The derivative expansion used for even numerators bounds the degree by
$\sum_{i=0}^{a-1}2i-\binom a2=\binom a2$.
For $a=0$ the minor is empty and has degree zero.
Because $L>\binom a2$, finite differences prove polynomiality of $Q$.
Its coefficients are integers directly from its defining series.

Suppose $d\ge1$. The reflection $O_{-u-1}=-O_u$ gives
\[
 F(-d-N)=\eta F(N),\qquad \eta=(-1)^{\binom d2}.
\]
In fact the columns reverse and each acquires a minus sign. The latter
$d$ signs cancel the change $(-1)^{d^2}$ of the external prefactor.
For $1\le s<d$ both $O_{-1}$ and $O_0$ occur in $F(-s)$, so the
determinant vanishes. Bilateral rational reciprocity, proved above,
therefore yields
\[
 G(1/z)=-\sum_{s\ge1}F(-s)z^s=-\eta z^dG(z).
\]
Multiplying by the factors defining $Q$ gives
$z^{pL-d}Q(1/z)=-\eta(-\varepsilon)^LQ(z)$.
As $Q(0)=1$, this identity also proves that the coefficient at $pL-d$
is the displayed sign and that no higher coefficient occurs.
The polygonal expression follows by substitution.
Reducing the sign exponent modulo two in the sixteen classes
$(k,m)\bmod4$ gives the stated list. Finally $d=0$ forces $t=a=0$,
so $k=1,m=0$. The moment Gram determinant then equals one at every
order. Hence $G(z)=(1-z)^{-1}$ and $Q(z)=1$, as claimed.


\begin{corollary}
For every $a=m+k-1\ge0$, the odd numerator in Cigler's normalization satisfies
\[
             (1-\varepsilon z^{2k-1})^{a}\mid Q(z)
             \qquad\text{in }\mathbb Z[z].
\]
\end{corollary}
\begin{proof}
The shift-polynomial degree bound $\binom a2$ in the preceding proof shows that $(1-\varepsilon z^p)^{1+\binom a2}G(z)$ is an integer polynomial. The difference between this exponent and the exponent $L$ defining $Q$ is exactly $a$.
\end{proof}

\section{A finite-difference obstruction for a binomial determinant}
Cigler's Conjecture~18 (arXiv:2308.07642v2, p.~12) also concerns
the sequence $c_{r,h}=\binom{2h+r}{h}$ for $h\ge0$, with $c_{r,h}=0$
for $h<0$ by the ordinary convention for a negative lower binomial index.
Write $d_{r,s}(N)=\det(c_{r,i+j+s})_{0\le i,j<N}$.
In the odd case $k=2$, $m=1$, its assertion specializes to the claim
that $q(n)=d_{7,-1}(7n+3)$ is a cubic with leading coefficient
\[
 (-1)^{\binom12+1}
 \frac{2^{3\cdot2-2}(2\cdot 2+1+1)^{2\cdot2-2}B_2}
 {(2\cdot2-2)!}=-48,
 \qquad B_2=\frac16.
\]
The printed factor in this coefficient is $(2k+m+1)^{2k-2}$.
The claim would give $\Delta^3q(0)=-288\equiv2\pmod5$.

The following elimination certificate uses arithmetic in $\mathbb F_5$.
In column $i$ take the first available row at or below $i$ with
nonzero entry, swap it with row $i$ if necessary and eliminate below
the pivot. A terminal zero records a column with no pivot.
The product of pivots times the row-swap sign gives the determinant.
\[
\begin{array}{r|l|r|r}
N&\text{successive pivots}&\text{swaps}& d_{7,-1}(N)\bmod 5\\\hline
3  &\mathtt{110}&1&0\\
10 &\mathtt{1111411114}&3&4\\
17 &\mathtt{1111411114140}&3&0\\
24 &\mathtt{111141111414111111414111}&7&1
\end{array}
\]
For a direct reproduction, the following standard-library Python
performs the same operations.
\begin{verbatim}
from math import comb
def b(h):
    return comb(2*h+7, h) % 5 if h >= 0 else 0
def det5(N):
    A = [[b(i+j-1) for j in range(N)] for i in range(N)]
    value = 1
    for i in range(N):
        j = next((r for r in range(i,N) if A[r][i]), None)
        if j is None: return 0
        if j != i:
            A[i], A[j] = A[j], A[i]
            value = -value
        pivot = A[i][i]
        value = value * pivot % 5
        for r in range(i+1,N):
            c = A[r][i] * pow(pivot,-1,5) % 5
            A[r] = [(x-c*y) % 5 for x,y in zip(A[r],A[i])]
    return value
assert [det5(N) for N in (3,10,17,24)] == [0,4,0,1]
\end{verbatim}
It follows that $\Delta^3q(0)\equiv1-3\cdot0+3\cdot4-0
\equiv3\pmod5$. This contradicts a necessary identity for the stated
cubic and refutes Conjecture~18 as printed. The even-order formula and
the problem of finding a general replacement for the odd coefficient
are separate questions.


\section{Complementary binomial rows and Sun's Conjecture 35(i)}
For $x\in\mathbb Z$ and $j\ge0$ use the falling-factorial definition of $\binom{x}{j}$, also for negative $x$. A rational polynomial is integer-valued when it takes integer values at every integer. In the variables of Sun's printed paper \cite[Conjecture 35(i)]{Sun}, the statement is as follows.
\begin{theorem}\label{thm:sun}
For every $\varepsilon\in\{1,-1\}$ and every $\ell,m,n\ge1$, the polynomial
\[
 \frac1n\sum_{k=0}^{n-1}\varepsilon^k(2k+1)^{2\ell-1}
       \sum_{j=0}^k\binom{-x}{j}^{m}\binom{x-1}{k-j}^{m}
\]
is integer-valued.
\end{theorem}
Guo \cite{Guo} uses the translated variable $x+1$ in Sun's expression, explains the Vandermonde case $m=1$, and proves $m=2$. Xia \cite{Xia} treats the distinct part 35(ii), with squares and a different modulus. Guo and Wang \cite{GW} already use the triangular phase below in single-index products. Our argument uses this phase on the sum of two independently reflected indices.

\begin{problem}
For two nonnegative upper entries whose sum is $-1$ modulo $n$, find a weight on a triangular Gaussian-binomial convolution that forces divisibility by $[n]_q$ and allows independently chosen powers on its two rows.
\end{problem}

Write $[s]_q=1+q+\dots+q^{s-1}$ for $s\ge1$, and use the Gaussian polynomials $\qbin Au$, set to zero when $u>A$. Let $T(k)=k(k+1)/2$ and put $\delta=0$ for $\varepsilon=1$, $\delta=1$ for $\varepsilon=-1$.
\begin{theorem}\label{thm:q}
Suppose $n\ge1$, $A,B\ge0$, $n\mid A+B+1$, and $h$ is a positive odd integer. For every polynomial $P(X,Y)\in XY\mathbb Z[X,Y]$, the Laurent polynomial
\[
 S(q)=\sum_{\substack{u,v\ge0\\u+v<n}}
 \varepsilon^k q^{\delta T(k)-hk}[2k+1]_q^h
       P\!\left(\qbin Au,\qbin Bv\right),\qquad k=u+v,
\]
belongs to $[n]_q\mathbb Z[q,q^{-1}]$.
\end{theorem}

\begin{proof}
The assertion is immediate for $n=1$. For $d>1$ dividing $n$, fix a primitive $d$th root $\zeta$. We use the Gaussian Lucas evaluation
\[
 \left.\qbin{ad+b}{cd+e}\right|_{q=\zeta}
 =\binom ac\left.\qbin be\right|_{q=\zeta}
 \quad(0\le b,e<d).
\]
For completeness, expand $\prod_{i=0}^{ad+b-1}(1+zq^i)$ by the finite $q$-binomial theorem. At $q=\zeta$ each full block gives $1-(-z)^d$. The coefficient of $z^{cd+e}$ is the product of the high binomial, the low Gaussian and its phase. Indeed
$\zeta^{\binom{cd+e}{2}-\binom e2}=(-1)^{c(d-1)}$:
for odd $d$ it is one, and for even $d$ it is $(-1)^c$. This exactly cancels the phase from the full blocks. The finite product expansion itself follows by multiplying the factors one at a time.

Write $A=ad+\alpha$, $B=bd+\beta$, $0\le\alpha,\beta<d$; then $\alpha+\beta=d-1$. Set $u=id+r$, $v=jd+s$. A nonzero summand must have $r\le\alpha$ and $s\le\beta$, because $P$ vanishes on both axes. The cutoff $u+v<n$ now says just $i+j<n/d$. Every remaining block contains the entire rectangle of low digits, including any entries on which $P$ happens to vanish. Reflect within that rectangle by
\[
    (r,s)\longmapsto(\alpha-r,\beta-s).
\]
The two Gaussian values are individually invariant, by the Lucas evaluation and the symmetry of each low Gaussian polynomial. Write $k'$ for the reflected total. With $H=i+j$ one has
\[
 k+k'+1=(2H+1)d,\qquad k'-k\equiv d-1\pmod 2.
\]
Consequently $[2k'+1]_\zeta=-\zeta^{-(2k+1)}[2k+1]_\zeta$, so that
\[
 \zeta^{-hk'}[2k'+1]_\zeta^h=-\zeta^{-hk}[2k+1]_\zeta^h.
\]
The alternating factor, when present, is invariant. Indeed $T(k')-T(k)=(k'-k)(2H+1)d/2$. For odd $d$ the difference is a multiple of $d$ and the two parities agree. For even $d$ the two parities differ and the extra power of $\zeta$ is $-1$. Thus each pair cancels, and any fixed term vanishes. This proves $S(\zeta)=0$.

Clear negative powers by a monomial. All factors $\Phi_d(q)$ for $d\mid n$ and $d>1$ divide the resulting integer polynomial. They are distinct monic irreducibles, so their product $[n]_q$ divides it in $\mathbb Z[q]$. Removing the monomial proves the assertion.
\end{proof}

\begin{proof}[Proof of Theorem~\ref{thm:sun}]
For a fixed integer $x$, take nonnegative representatives $A\equiv-x$, $B\equiv x-1\pmod{n!}$. For every $0\le j<n$, the integer polynomial $X(X-1)\cdots(X-j+1)$ preserves congruences modulo $n!$. Since $nj!\mid n!$, division by $j!$ gives
\[
 \binom Aj\equiv\binom{-x}j\pmod n,\qquad
 \binom Bj\equiv\binom{x-1}j\pmod n .
\]
Apply Theorem~\ref{thm:q} with $P(X,Y)=X^mY^m$ and $h=2\ell-1$, and evaluate at $q=1$. The Gaussian entries specialize to ordinary binomial coefficients and $[n]_1=n$. Regrouping the pairs by their total recovers Sun's exact numerator.
\end{proof}

The same argument applies with independent exponents $m_1,m_2\ge1$ and any odd polynomial weight in $2k+1$ with integer coefficients. Both small features of the proof matter. For $P(X,Y)=XY$ at $n=2,A=0,B=1,h=1$ and $\varepsilon=-1$, omitting $q^{T(k)}$ leaves value $2$ at $q=-1$. Also the modulus $n$ is best possible uniformly: at $x=0$, $h=1$ and $\varepsilon=(-1)^{m+1}$ the numerator is $\sum_{k=0}^{n-1}(-1)^k(2k+1)=(-1)^{n-1}n$.


\subsection{A squarefree multiplier in the square case}
Xia \cite{Xia} proved Sun 35(ii) and Guo's coefficient conjecture, and obtained additional binomial factors in the modulus with the original multiplier. The following elementary argument gives a squarefree multiplier at modulus $n^2$.
For $\ell\ge1$ put $P_\ell=\prod_{3\le p\le 2\ell-1,\ p\text{ prime}}p$,
with empty product one. Write $C_k=\binom{2k}{k}$.
\begin{theorem}
For $n,\ell\ge1$ and $0\le k<n$,
\[
 n^2\ \bigm|\ P_\ell C_k^2
 \sum_{m=k}^{n-1}(2m+1)^{2\ell-1}\binom{m+k}{2k}.
\]
Consequently Guo's Conjecture 2, equation (4.1) of his paper on integer-valued
polynomials, and Sun's Conjecture 35(ii) hold. In each, $P_\ell$ suffices in
place of the stated odd double factorial.
\end{theorem}
\begin{proof}
Put $Q_j(m)=\binom{m+j}{2j}$ for nonnegative integers $m,j$,
zero when $j>m$. Two identities will be used:
\begin{align*}
 (2m+1)^2Q_j(m)&=(2j+1)^2Q_j(m)+4(2j+1)(2j+2)Q_{j+1}(m),\\
 \frac{1}{n^2}\sum_{m=0}^{n-1}(2m+1)Q_j(m)
 &=\frac{\binom{n-1}{j}\binom{n+j}{j}}{(j+1)C_j}\qquad(j<n).
\end{align*}
The first follows by subtracting $(2j+1)^2$. For the second,
$U_j(m)=m(m-j)Q_j(m)/(j+1)$ has first difference
$U_j(m+1)-U_j(m)=(2m+1)Q_j(m)$, including when $m<j$.
Summing gives the stated factorization. For $j\ge n$ the sum vanishes.

After $\ell-1$ applications of the first identity,
the coefficient of $Q_{k+r}(m)$ in $(2m+1)^{2\ell-2}Q_k(m)$
is an integer multiple of $4^r(2k+2r)!/(2k)!$.
Indeed each contribution raises the index at exactly the levels $k,\ldots,k+r-1$,
while the steps at a fixed level have integral weights. Here $0\le r\le\ell-1$.
In the second identity put $j=k+r$. After cancellation it suffices to prove
\[
 k+r+1\ \bigm|\ P_\ell 4^r C_k \big((k+r)!/k!\big)^2 .
\]

Fix a prime $p$ and let $p^e\mid k+r+1$. Put
$h=\lfloor\log_p(2r+1)\rfloor$. For $h<i\le e$, the least residue
of $k$ modulo $p^i$ is $p^i-1-r\ge p^i/2$, since $p^i\ge2r+2$.
Legendre's factorial formula therefore contributes one to
$v_p(C_k)$ at each such $i$. Hence $v_p(C_k)\ge\max(0,e-h)$.
For odd $p$ and $h\ge1$ we have $p\le2r+1\le2\ell-1$ and
$p^{h-1}\le r$. The product $(k+r)!/k!=r!\binom{k+r}{r}$ is divisible
by $p^{h-1}$, and $P_\ell$ supplies a further factor $p$.
For $p=2$, $h\le2r$ because $2r+1\le 2^{2r}$, so $4^r$ covers
the deficit. This proves the divisibility.

Guo's transformation, his equation (2.1), with his variable set to $x-1$,
gives the exact identity
\[
 \sum_{j=0}^{m}\binom{-x}{j}^2\binom{x-1}{m-j}^2
 =\sum_{k=0}^{m}Q_k(m) C_k^2\binom{x+k-1}{2k}.
\]
For every integral $x$ the last factor is integral, also at negative tops.
Summing the coefficient congruences proves the strengthened form of Sun
35(ii). Finally $P_\ell\mid (2\ell-1)!!$.
\end{proof}


\subsection{A Gaussian analogue for every odd weight}
After equation (4.2), Guo asks for a $q$-analogue of his coefficient congruence
in higher odd degrees. A factor depending on the exponent in the binomial
expansion of the weight deals with roots of unity of even order.

Work in $\mathbb Z[v,v^{-1}]$, put $q=v^2$, and use symmetric integers
$s_b=(v^b-v^{-b})/(v-v^{-1})$. Their factorials are denoted by $s_b!$.
Write $B(u,w)=s_u!/(s_w!s_{u-w}!)$ for $0\le w\le u$ and zero otherwise.
These are ordinary Gaussian coefficients in $q$ up to a monomial in $v$.
For $\ell\ge1$ put
\[
 X_m=s_ms_{m+1},\qquad
 W_{\ell,m}=\sum_{a=0}^{\ell-1} B(\ell-1,a)
       \left(\prod_{h=1}^{a}(1+q^h)\right)^2 X_m^a .
\]
At $v=1$, the expression $W_{\ell,m}$ is
$(1+4m(m+1))^{\ell-1}=(2m+1)^{2\ell-2}$.

\begin{theorem}
For $\ell,n\ge1$ and $0\le k<n$, the Laurent polynomial
\[
 \left(\prod_{h=1}^{\ell}s_{2h-1}\right) B(2k,k)^2
     \sum_{m=k}^{n-1}s_{2m+1}W_{\ell,m}\,B(m+k,2k)
\]
is divisible by $s_n^2$. At $\ell=1$ this is Guo's (4.2) up to a
Laurent monomial, and specialization at $v=1$ gives his exact (4.1).
\end{theorem}
\begin{proof}
Let $C_j^*=B(2j,j)$ and $Q_j^*(m)=B(m+j,2j)$. The relations
\begin{align*}
 X_mQ_j^*(m)&=X_jQ_j^*(m)+s_{2j+1}s_{2j+2}Q_{j+1}^*(m),\\
 \sum_{m=0}^{n-1}s_{2m+1}Q_j^*(m)
 &=s_n^2\frac{B(n-1,j)B(n+j,j)}{s_{j+1}C_j^*}\qquad(j<n)
\end{align*}
follow respectively from $X_m-X_j=s_{m-j}s_{m+j+1}$ and telescoping
$s_m s_{m-j}Q_j^*(m)/s_{j+1}$. The latter uses
$s_{m+1}s_{m+j+1}-s_ms_{m-j}=s_{2m+1}s_{j+1}$.
In the expansion of $X_m^aQ_k^*(m)$, the coefficient of $Q_j^*(m)$,
$j=k+r$, has the factor $s_{2j}!/s_{2k}!$, since each index between
$k$ and $j-1$ must be raised once. After inserting the second identity,
the corresponding quotient by $s_n^2$ is a Laurent polynomial times
\[
 \frac{1}{s_{j+1}}
 \left(\prod_{h=1}^{\ell}s_{2h-1}\right)
 \left(\prod_{h=1}^{a}(1+q^h)\right)^2
 C_k^* \left(\frac{s_j!}{s_k!}\right)^2 .
\]

Up to a Laurent monomial, $s_{j+1}$ is the product of $\Phi_d(q)$
over $d>1$ dividing $j+1$. If $d>2r+1$, the least residue of
$k=j-r$ modulo $d$ is $d-1-r\ge d/2$. The exponent of $\Phi_d(q)$ in
the ordinary central Gaussian coefficient is
$\lfloor2k/d\rfloor-2\lfloor k/d\rfloor=1$.
If $d\le2r+1$ is odd, the numerator contains $s_d$ because
$2r+1\le2\ell-1$. If $d\le2r+1$ is even, then
$d/2\le r\le a$, so the numerator contains $1+q^{d/2}$.
Thus every distinct cyclotomic factor of the denominator is covered in
$\mathbb Z[q]$ after removing units, and the resulting divisibility also
holds under $q=v^2$.
\end{proof}


\section{The first two prime-square congruences in part 35(iii)}
\label{sec:cubicprime}
For a prime $p>3$ and $x\in\mathbb Z_p$, write
\[
 T_p(x)=\sum_{k=0}^{p-1}(-1)^k(2k+1)
        \sum_{j=0}^{k}\binom{-x}{j}^{3}\binom{x-1}{k-j}^{3}.
\]
The first two claims in Sun's Conjecture 35(iii) have the following joint evaluation.
Let the Legendre symbol be $\chi=\bigl(\frac{p}{3}\bigr)$ and $c=(p\chi-1)/3$. Then
\begin{equation}\label{eq:cubic-sun-eval}
 T_p(x)\equiv
 \begin{cases}
   0 & 3x\not\equiv 1,2\pmod p,\\
   x+c & 3x\equiv1\pmod p,\\
   1-x+c & 3x\equiv2\pmod p
 \end{cases}
 \qquad (\bmod\ p^2).
\end{equation}
The third line spells out the symmetry $T_p(x)=T_p(1-x)$.

\subsection*{Two one-variable evaluations}
For complex $z$ in a neighborhood of a nonnegative integer, set
\[
 F(z)=\sum_{j\geq0}(-1)^j\binom zj^3,\qquad
 K(z)=\sum_{j\geq0}(2j-z)(-1)^j\binom zj^3.
\]
The Dixon and Dougall summations give
\begin{align}
 F(z)&=\cos(\pi z/2)\frac{\Gamma(1+3z/2)}{\Gamma(1+z/2)^3},\label{eq:suncubic-dixon}\\
 K(z)&=-2\sin(\pi z/2)\frac{\Gamma((1+3z)/2)}{\Gamma((1+z)/2)^3}.\label{eq:suncubic-dougall}
\end{align}
Indeed, take $a=b=c=-z$ in Dixon's $\mathstrut{}_3F_2(1)$ summation.
In Dougall's $\mathstrut{}_5F_4(1)$ summation, with upper parameters
$a,1+a/2,b,c,d$, take $a=b=c=-z$ and $d=(1-z)/2$.
The last parameter cancels the lower parameter $1+a-d$. Multiplying the sum by $-z$ introduces the factor $2j-z$.
Euler's reflection formula gives the displayed expressions.
Both summation formulas are recorded together by Masjed-Jamei and Koepf
\cite[eqs.~(8), (13)]{masjedkoepf}.
The parameter substitutions hold off the isolated singularities of the hypergeometric notation and extend there by continuity.
For $\Re z>-1/3$ both series converge locally uniformly. Indeed,
$(-1)^j\binom zj=\Gamma(j-z)/(\Gamma(-z)\Gamma(j+1))$ is
$O(j^{-1-\Re z})$ locally uniformly for large $j$, by the gamma-ratio estimate.

For $N\ge0$ and $j>N$, $\binom zj^3$ has a zero of order three at $z=N$.
Consequently the values and first derivatives of $F,K$ at $N$ agree with those of their sums through $j=N$.
In particular these derivatives are rational numbers, obtained by differentiating finite polynomials.
The gamma formulas yield $K(2u)=F(2u+1)=0$ and
\begin{align*}
 F(2u)&=(-1)^u\frac{(3u)!}{u!^3},&
 K'(2u)&=-(-1)^u\frac{(6u)!\,u!^3}{(3u)!\,(2u)!^3},\\
 K(2u+1)&=-2(-1)^u\frac{(3u+1)!}{u!^3},&
 F'(2u+1)&=-2(-1)^u\frac{(6u+4)!\,(u+1)!^3}{(3u+2)!\,(2u+2)!^3}.
\end{align*}
These follow directly from $\Gamma(t+1/2)=(2t)!\sqrt\pi/(4^t t!)$ for integers $t\geq0$.

\subsection*{Arithmetic of complementary rows}
Put $X=-x$, $Y=x-1$, and take $A,B\in[0,p-1]$ representing $X,Y$
modulo $p$. Thus $A+B=p-1$ and $A,B$ have the same parity. Write
$X=A+p\sigma$, $Y=B+p\tau$ with $\sigma,\tau\in\mathbb Z_p$.
All binomials of degree less than $p$ have $p$-integral coefficients.
For $A<j<p$, $\binom Xj^3\equiv0\pmod{p^3}$, and the same applies to the second row beyond $B$.
The surviving rectangle $0\leq i\leq A$, $0\leq j\leq B$ lies in $i+j<p$.
Since $2(i+j)+1=(2i-X)+(2j-Y)$, its weighted double sum factors into two products,
the centered moment in either row times the zeroth moment in the other.
Taylor expanding the finite polynomials around $A,B$ gives
\[
 T_p(x)\equiv p\sigma\,C(A,B)+p\tau\,C(B,A)\pmod{p^2},
 \qquad
 C(A,B)=\begin{cases}K'(A)F(B)&A\text{ even},\\
 F'(A)K(B)&A\text{ odd}.
 \end{cases}
\]
The Taylor expansions take place in $\mathbb Z_p$ and use only the four rational formulas above.

Set $w=3A$ for even $A$, and $w=3A+1$ for odd $A$. Legendre's factorial formula gives
\[
 v_p(C(A,B))=
 \left\lfloor\frac w p\right\rfloor-
 \left\lfloor\frac w {2p}\right\rfloor+
 \left\lfloor\frac{3p-3-w}{2p}\right\rfloor .
\]
In fact the three factorials contributing here have arguments $w$, $w/2$ and $(3p-3-w)/2$.
All other factorial arguments are below $p$, and $w\leq3p-3<p^2$.
For $w\geq p$ the difference of the first two floors is at least one.
For $w<p$ the last floor vanishes precisely when $w>p-3$. Since $w$ is even, the only case with valuation zero is $w=p-1$.
It follows that both products are divisible by $p$ unless $3x$ is $1$ or $2$ modulo $p$.

It remains to evaluate the sole unit in each exceptional case. By symmetry take $3x\equiv1\pmod p$.
When $p=6u+1$ one has $A=2u$, $B=4u$ and
\[
 C(A,B)=-(-1)^{3u}\frac{(6u)!^2\,u!^3}{(3u)!\,(2u)!^6}.
\]
Splitting $(2u)!$ into even and odd factors and reducing each odd factor $2j-1$ as $-2(3u+1-j)$ gives
$ (2u)!^2\equiv(-1)^u2^{2u}u!(3u)!\pmod p$.
Wilson's theorem yields $(((p-1)/2)!)^4\equiv1\pmod p$ and $(p-1)!^2\equiv1\pmod p$.
Together with $2^{6u}\equiv1\pmod p$ these identities give $C(A,B)\equiv-1\pmod p$.
Hence $T_p(x)\equiv-p\sigma=x+(p-1)/3\pmod{p^2}$.

When $p=6u+5$ one has $A=4u+3$, $B=2u+1$. Cancelling $(u+1)^3/(2u+2)^3$ in the factorial expression gives
\[
 C(B,A)=\frac{(-1)^{3u+1}}{2}
          \frac{(6u+4)!^2\,u!^3}{(3u+2)!\,(2u+1)!^6}.
\]
Here the odd factors of $(2u+1)!$ reduce to $-2(3u+3-j)$ for $1\leq j\leq u+1$.
Thus $(2u+1)!^2\equiv(-1)^{u+1}2^{2u+1}u!(3u+2)!\pmod p$.
Wilson's theorem and $2^{6u+3}\equiv1/2\pmod p$ now give $C(B,A)\equiv1\pmod p$.
This yields $T_p(x)\equiv p\tau=x-(p+1)/3\pmod{p^2}$ and completes the proof of \eqref{eq:cubic-sun-eval}.


\subsection{The prime-square assertions in Sun's Conjecture 35(iv)}
Define $T_p^{[3]}(x)$ by replacing the factor $2k+1$ in $T_p(x)$ by $(2k+1)^3$. We prove
\[
 T_p^{[3]}(x)\equiv0\pmod{p^2}
 \quad\text{if }3x\not\equiv -1,1,2,4\pmod p
\]
for every odd prime $p$. We also prove the same conclusion for $p>5$ and $x=\pm(p-(\frac p3))/3$. In the proof we write $a=x$.
For an odd prime $p$ and $a\in\mathbb Z_p$, write
$M_e(a)=\sum_{j=0}^{p-1}(2j+a)^e(a)_j^3/(j!)^3$.
Here $(a)_j$ is the rising factorial.
For the first assertion of part (iv) take $b=1-a$ and
$3a\not\equiv-1,1,2,4\pmod p$.
The separate product vanishings from the preceding calculation of part (iii) give
\[
 M_1(a)M_0(b)\equiv M_0(a)M_1(b)\equiv0\pmod{p^2}\quad(p>3).
\]
The cube of the weight can be reduced to these same two products.
Indeed, for $c_j=(a)_j^3/(j!)^3$, and $d=0,1$, summation of
$j^3 c_j=(j-1+a)^3c_{j-1}$ gives
\[
 \sum_{j=0}^{p-1}\big(j^{d+3}-(j+1)^d(j+a)^3\big)c_j
 =-p^d(p-1+a)^3c_{p-1}\equiv0\pmod{p^3}.
\]
The numerator of $(a)_p$ contains a multiple of $p$. In terms of $2j+a$,
the $d=0$ identity gives
$M_2(a)\equiv-a^2M_0(a)/3\pmod{p^2}$ when $a$ is a unit.
When $p\mid a$ the same formula follows directly, because the terms
$c_j$ for $j\ge1$ are divisible by $p^3$. Substituting this expression
in the $d=1$ identity gives
\[
 M_3(a)\equiv-\frac{a^2(a+3)}{3a+1}M_1(a)\pmod{p^2}.
\]
Both formulas also apply with $a$ replaced by $b$ under the stated exclusions.
With $w=2(i+j)+1$, the original summand equals $w^3c_i(a)c_j(b)$.
For $i,j<p$ with $i+j\ge p$, at least one of $i,j$ exceeds the first
nonnegative residue of $-a,-b$, respectively, since those residues sum to
$p-1$. Every discarded term is therefore divisible by $p^3$.
Expanding $w^3=((2i+a)+(2j+b))^3$ on the resulting rectangular range
and applying the two reductions yields
\begin{align*}
 &-\left(\frac{a^2(a+3)}{3a+1}+b^2\right)M_1(a)M_0(b)\\
 &\hspace{7mm}-\left(\frac{b^2(b+3)}{3b+1}+a^2\right)M_0(a)M_1(b)
 \equiv0\pmod{p^2}.
\end{align*}
For $p=3$, direct substitution of the nine residues for $a$ modulo nine in
the six terms $i+j<3$ gives zero in every case. Equivalently, modulo nine
the terms with $i+j=1$ vanish by their cubic weight; the remaining sum is
$1+5^3((a(a+1)/2)^3+(ab)^3+(b(b+1)/2)^3)$, with $b=1-a$.
Cubes of integers modulo nine depend only on their residue modulo three,
so verification needs only $a=0,1,2$. This finishes the generic cubic
assertion of part (iv). We next prove the two specified integral exceptions.

The final pair of exceptional integral arguments in part (iv) also satisfies
the stated congruence. Let $p>5$, $\chi=\bigl(\frac{p}{3}\bigr)\in\{-1,1\}$ and
$h=(p-\chi)/3$, an even integer. Of the choices $a=\pm h$, the one with
$3a\equiv1\pmod p$ is $a=(1-\chi p)/3$. The residue-class calculation
in part (iii) shows that at most one of the two separate first-moment
products can contribute modulo $p^2$, and their sum there is
$a+(p\chi-1)/3=0$. They both vanish. The two cubic reductions have unit
denominators and give the same conclusion for the cubic convolution.

For the choice $3a\equiv-1\pmod p$, the separate first-moment products
vanish by the generic case of part (iii). The denominator $3b+1$ is a unit
because it is $5$ modulo $p$. The sole term for which the displayed cubic
reduction may fail is $M_3(a)M_0(b)$, and it vanishes identically.
If $p\equiv1\pmod3$, then $a=h$ and $b=-(h-1)$. For any odd integer
$0\le m<p$, the terms $j$ and $m-j$ in
$M_0(-m)=\sum_{j=0}^m(-1)^j\binom{m}{j}^3$ cancel, hence $M_0(b)=0$.
If $p\equiv2\pmod3$, then $a=-h$. For even $h$, the same reflection
preserves $(-1)^j\binom hj^3$ and negates $(2j-h)^3$, hence $M_3(a)=0$.
The remaining $a=1/2$ assertion modulo $p^5$ requires additional information.


\section{Examples, checks and the limits of the historical claim}

For $r=7$, Theorem~\ref{result:odd} gives degree three at determinant sizes $7n+2$. The first four signed values are $-14,315,1673,4746$. They determine the entire sequence,
\[
 (-1)^nD_{7,0}(7n+2)
 =-14+329n+1029\binom n2+686\binom n3.
\]
Here the degree theorem is essential: four values alone leave infinitely many possible continuations. For an order beyond the finite ranges discussed by Wang and Xin, the same theorem gives degree $155$ for both $D_{43,0}(43n+6)$ and $D_{43,0}(43n+38)$ after multiplication by $(-1)^n$. Thus 156 values determine each polynomial even though its determinantal definition has unbounded size.

For a check independent of the reductions, direct integer elimination on the original coefficient matrices verified the required numerator degrees and reciprocity for $1\le r\le7$ at the three shifts corresponding to $a=0,1,2$. The finite coefficient comparisons extended three or more places beyond each predicted degree. Direct integer summation also verified 7,020 unequal-power or equal-power congruences with $1\le n\le10$, $-6\le x\le6$, $1\le m_1,m_2\le3$, $h\in\{1,3,5\}$ and both signs. All deductions over unbounded parameter sets above are analytic or algebraic.

\paragraph{G1 and G2: targets and fidelity.}
The theorem numbers and source versions are identified in Section~1. The Hankel arguments start with the original coefficient determinant, include the empty determinant, and account for every initial zero for negative shifts. The proof of Sun 35(i) starts and ends with the printed complementary binomial entries. Our higher-degree Gaussian analogue responds to Guo's question after his equation (4.2). Our odd-numerator classification retains Cigler's original multiplier and identifies exactly where his predicted degree fails.

\paragraph{G3: public status.}
Wang and Xin established bounded ranges and described further computations; Cigler subsequently stated the unrestricted degree and polynomiality questions. Guo proved the square case of Sun 35(i). In the original publications, forward citation records, related preprints, problem dossiers and corpus texts examined for this manuscript, we found no earlier full solution of Wang--Xin 9 or Sun 35(i). This records the scope of a literature check through 20 November 2026. Xia's theorem already proves 35(ii), and the published counterexample to Cigler 16 already settles the truth value of its universal conjunction.

\paragraph{G4 and G5: proofs and range.}
The determinant proof supplies a nonzero leading coefficient in every requested residue, including cases in which signed permutation terms might cancel. The four exact determinant values in the binomial-determinant test give a counterexample to the universal conjunction in Cigler 18. The arithmetic proof covers all divisors of a composite modulus and both parities of the order of each root of unity. The prime-square statements in the last sections have their own proofs; the higher prime-power clauses in 35(iii) and 35(iv) are outside those deductions.

\paragraph{G6: provenance.}
Christoffel determinant reduction, Gaussian Lucas evaluation, positive moment matrices, Dixon's and Dougall's sums, and the triangular phase all have antecedents. The contribution made here is the identification and nonvanishing of the folded leading minor in both parities, the classification through bilateral reflection, and the complementary-digit reflection of a full two-dimensional convolution. Historical priority for the results and for these combinations remains open to bibliographic correction. This manuscript has not undergone external peer review.

\begingroup\small
\begin{thebibliography}{99}
\setlength{\itemsep}{2pt}\setlength{\parsep}{0pt}
\bibitem{WX} Y. Wang and G. Xin, Hankel determinants for convolution powers of Catalan numbers, \emph{Discrete Math.} \textbf{342} (2019), 2694--2716. \href{https://doi.org/10.1016/j.disc.2019.05.025}{doi:10.1016/j.disc.2019.05.025}. \href{https://arxiv.org/abs/1811.00248}{arXiv:1811.00248}.
\bibitem{Cig} J. Cigler, Some experimental observations about Hankel determinants of convolution powers of Catalan numbers, version 2, August 2023. \href{https://arxiv.org/abs/2308.07642v2}{arXiv:2308.07642v2}.
\bibitem{Back} J. Cigler, Hankel determinants of convolution powers of Catalan numbers revisited, 2024. \href{https://arxiv.org/abs/2403.11244}{arXiv:2403.11244}. See also M. Fulmek's bijective treatment, \href{https://arxiv.org/abs/2402.19127}{arXiv:2402.19127}.
\bibitem{Krat} C. Krattenthaler, Hankel determinants of linear combinations of moments of orthogonal polynomials, II, \emph{Ramanujan J.} \textbf{61} (2023), 597--627. \href{https://arxiv.org/abs/2101.04225}{arXiv:2101.04225}.
\bibitem{Other} S. Chern and W. Shi, Hankel determinants of Catalan-like sequences, 2026. \href{https://arxiv.org/abs/2608.27208}{arXiv:2608.27208}. G.-N. Han, Hankel determinants for convolution powers of Narayana polynomials, \emph{Electron. J. Combin.} \textbf{33} (2026), P2.42. \href{https://doi.org/10.37236/14995}{doi:10.37236/14995}.
\bibitem{Omega} The Omega Institute, public problem dossiers on \href{https://www.trureturing.org/research/cigler-catalan-power-shifted-hankel/}{Cigler's Conjecture 11} and \href{https://www.trureturing.org/research/cigler-sixteen-literal-numerator-degree/}{the literal degree in Conjecture 16}; see also the \href{https://the-omega-institute.github.io/trureturing-mdbook/Blueprint/D5/S3/Combinatorics/CatalanPowerHankel/CiglerEleven.html}{determinant-reduction blueprint}.
\bibitem{Sun} Z.-W. Sun, Open conjectures on congruences, \emph{Nanjing Univ. J. Math. Biquarterly} \textbf{36} (2019), 1--99, Conjecture 35, pp.~29--31. \href{https://doi.org/10.3969/j.issn.0469-5097.2019.01.01}{doi:10.3969/j.issn.0469-5097.2019.01.01}.
\bibitem{Guo} V. J. W. Guo, On a conjecture related to integer-valued polynomials, \emph{Bull. Math. Soc. Sci. Math. Roumanie} \textbf{66} (2023), 449--453. \href{https://ssmr.ro/bulletin/pdf/66-4/articol_8.pdf}{Publisher's text}.
\bibitem{GW} V. J. W. Guo and S.-D. Wang, Factors of sums and alternating sums of products of $q$-binomial coefficients and powers of $q$-integers. \href{https://doi.org/10.11650/tjm/180601}{doi:10.11650/tjm/180601}. \href{https://arxiv.org/abs/1705.06236}{arXiv:1705.06236}.
\bibitem{Xia} W. Xia, Proof of a conjecture of Sun and its extension by Guo, \emph{Ramanujan J.} \textbf{62} (2023), 617--631. \href{https://doi.org/10.1007/s11139-022-00668-z}{doi:10.1007/s11139-022-00668-z}.
\bibitem{masjedkoepf} M. Masjed-Jamei and W. Koepf, Some summation theorems for generalized hypergeometric functions, \emph{Axioms} \textbf{7} (2018), 38. \href{https://doi.org/10.3390/axioms7020038}{doi:10.3390/axioms7020038}.
\bibitem{Corpus} QEEDDD, \emph{Verifiable Mathematical Reasoning Corpus}, \href{https://github.com/QEEDDD/Verifiable-Mathematical-Reasoning-Corpus/tree/ac1492a50784be919fe0b5e8f19f96f8d221a22e}{snapshot ac1492a}, 20 November 2026. See manuscripts 026, 079 and 208.
\end{thebibliography}\endgroup
\end{document}
